% =====================================================================
%  Chapter 14 — Indirect Detection and the Anomaly Record
% =====================================================================

\chapter{Indirect Detection and the Anomaly Record}
\label{ch:indirect}

The indirect channel looks for the products of dark matter annihilation or
decay: photons, neutrinos, charged particles. It has the highest
signal-to-noise ratio of any method in the field, because the ratio of
signal to background in the gamma-ray sky near the galactic centre can be
of order a percent. It also has the highest rate of non-replication in the
field's history, and the reason is instructive.

\section{How the channel works}If dark matter annihilates into Standard Model
particles at a rate $\langle\sigma v\rangle$ and a fraction of those
produce gamma rays, then the gamma-ray flux from a halo is

\[
\Phi(\varepsilon) \;=\; \frac{1}{2}\,
\frac{\langle \sigma v\rangle}{m_\chi^2}\,
\frac{dN_\gamma}{d\varepsilon}\,
\int_{E_{\rm min}}^{E_{\rm max}} \frac{\rho_\chi(r)}{r^2}\,dV .
\]

The integral is the J-factor, which depends on the halo density profile and
on the line of sight. Everything on the right except $\langle\sigma v\rangle$
and $m_\chi$ is astrophysics, and the astrophysics is where the difficulty
lives.

\section{The galactic centre excess}The most significant indirect result of the last
decade is the persistent excess of gamma rays above about 10~GeV in the
direction of the galactic centre, first identified in the Fermi-LAT data and
since confirmed with better statistics.

Totani's analysis, published in the Journal of Cosmology and Astroparticle
Physics in November 2025, fits the excess with a dark matter annihilation
model with a particle mass of 0.5 to 0.8~TeV and an effective annihilation
cross section of roughly 5 to $8\times 10^{-25}\ \cmthree\s$
\idx{Totani}. The statistical quality of the fit is good, the analysis is
careful, and it is a better piece of work than most of its genre.

It is also, in its central interpretation, undercut. The AMS-02 antiproton
flux constrains the annihilation cross section at roughly forty times lower
than the value the gamma-ray fit requires, with the measured antiproton
spectrum deviating from the Standard Solar Model by more than 20 sigma in
the relevant energy range \idx{Antiproton constraints}. The antiproton
measurement and the gamma-ray measurement are independent, and they disagree
about the dark matter interpretation. The gamma-ray excess itself is not in
doubt; the dark matter interpretation is.

This is the pattern to understand. The measurement is real; the modelling is
reasonable; the interpretation is over-extended.

\section{The positron excess}The rise in the positron fraction above 10~GeV in the
AMS-02 data is, by common assessment, established \idx{Positron excess}. What
is contested is the explanation: pulsars, nearby supernova remnants, and
dark matter annihilation all produce positrons, and the shape of the spectrum
above 500~GeV and the absence of a corresponding antiproton excess are the
discriminating information, and they do not fit the annihilation template
well.

The positrons are a case where the correct statement is that the measurement
is settled and the interpretation is not. It is also a case where the
ambiguity is irreducible: the sources are astrophysical objects that the
galaxy is full of, and no improvement in the dark matter model will resolve
it. Only a measurement that can separate the components --- and the
gamma-ray sky is the most promising route --- will.

\section{The 3.5 keV line}The claim that an X-ray emission line at 3.5~keV
detected in the spectra of galaxy clusters was the signature of decaying
dark matter became one of the most cited results of the 2010s. It was not
confirmed. A re-analysis of the same clusters with a different stacking
method found no line \idx{3.5 keV non-confirmation}, and a dedicated
observation of the Draco dwarf galaxy found a line at the same energy in one
galaxy and not in another, with a signal-to-noise ratio inconsistent with
the cosmological origin the claim required \idx{Draco line}.

The 3.5 keV case is instructive in a specific way: the original claim was
made by stacking 15 galaxy clusters, and the non-confirmation used the same
data with a different analysis. Nothing about the universe changed between
the two papers.

\section{The 21 cm anomaly and EDGES}In 2018 two experiments reported the
flattening of the 21~cm absorption profile at redshifts around 17, one of
them at a level of 0.1 per cent of the expected absorption, a signal that
would be produced by dark matter scattering on hydrogen at a cross section
of about $2\times 10^{-15}\ \cmtwo$. The SARAS 3 experiment's own follow-up
did not confirm its original result \idx{SARAS 3}, and the EDGES result was
strongly challenged on the grounds that the instrument's calibration and
foreground model could not support a 0.1 per cent effect
\idx{EDGES concerns}.

There is a further confusion in this part of the literature that should be
stated, because it has propagated: the 21~cm constraint most often cited
alongside the EDGES and SARAS claims is by Berlin, Hooper and Krnjaic,
Physical Review Letters 121, 011102 (2018) \idx{21 cm constraint}, and it is
frequently attributed to other authors.

\section{What the record teaches}The anomalies of the last five years share a structure, and
the structure is the lesson.

\begin{itemize}
\item Each was a real measurement of a real signal, published by careful
groups.
\item Each involved a modelling step that was reasonable but not
over-determined --- the halo profile, the foreground model, the
instrumental calibration.
\item Each was a claim about the entire sky, and the look-elsewhere effect
was not fully controlled.
\item None replicated, and in two cases the non-confirmation used the same
data as the original claim.
\end{itemize}

Nothing in this record suggests misconduct. The errors were statistical and
instrumental, and they were errors that the field's own culture of
speculation made more likely. The lesson for the 248~keV candidate is
direct: a single event in an unblinded window is the pattern, and the
pattern has a base rate.

\begin{keytakeaways}
\item The indirect channel has the best signal-to-noise in the field and the
worst replication record, because it depends on astrophysical modelling over
the whole sky.
\item The Fermi-LAT galactic centre excess is a good fit that the
antiproton measurement excludes as a dark matter interpretation.
\item The 3.5 keV line, the EDGES signal, the 21 cm anomaly and the SARAS 3
result were all real measurements that did not replicate, none of them
through misconduct.
\item The look-elsewhere effect and the foreground model were the recurring
failures, and the response is to require replication before interpretation.
\end{keytakeaways}
